%%%%%%%%%%%%%%%%%%%%%%%
%%% DOCUMENT SETTINGS
%%%%%%%%%%%%%%%%%%%%%%%
\documentclass[11pt]{exam}
%\printanswers % Uncomment to show answers
%%%%%%%%%%%%%%
%%% PACKAGES
%%%%%%%%%%%%%%
\usepackage{amsmath,amsfonts,amssymb,amsthm} % math fonts and symbols
\usepackage{bbm} % Math blackboard font
\usepackage{enumitem} % For reference enumerations lists
\usepackage{textcomp} % some symbols (like copyright)
\usepackage[top = 2cm, bottom = 2cm, right = 1.2cm, left = 1.2cm, paper = a4paper, headheight = 6cm]{geometry} % Margins setup
\usepackage{xcolor} % Colors
\usepackage{graphicx}
%%%%%%%%%%%%%%%%%%%%%%%%%%
%%% MACROS AND COMMANDS
%%%%%%%%%%%%%%%%%%%%%%%%%%
\newcommand{\N}{\mathbb{N}} % natural number set
\newcommand{\R}{\mathbb{R}} % real number set
\makeatletter
\renewcommand{\@makefnmark}{\textsuperscript{\arabic{footnote}}} % Ensure numeric superscripts
\renewcommand{\thefootnote}{\arabic{footnote}} % Ensure numeric labels
\makeatother
\setcounter{footnote}{1}
\newcommand{\BAR}{%
    \hspace{-\arraycolsep}%
    \strut\vrule % the `\vrule` is as high and deep as a strut
    \hspace{-\arraycolsep}%
}
%%%%%%%%%%%%%%%%%%%%%%%%%%%%%%%%%
%%% HEADER AND FOOT 
%%%%%%%%%%%%%%%%%%%%%%%%%%%%%%%%%
% Defining information
\newcommand{\degree}{Bachelor in Philosophy, Politics and Economics}
\newcommand{\shortdeg}{PPE}
\newcommand{\exam}{1st Midterm}
\newcommand{\subject}{Mathematics}
\newcommand{\theyear}{2025/26}
\newcommand{\ccauthor}{A. G. Azze - CUNEF Universidad}
\newcommand{\thedate}{Oct 10, 2025}
\newcommand{\university}{CUNEF Universidad}
% Setting header and footer
\pagestyle{headandfoot}
\header{\degree \\ \subject\ - \theyear}{}{\thedate \\ \ccauthor}
\footer{\university}{}{\thepage}
%%%%%%%%%%%%%%%%%%%%%%%%%
%%% DOCUMENT CONTENT
%%%%%%%%%%%%%%%%%%%%%%%%%
\begin{document}
%--- Instructions ---%
\begin{center}

    {\Huge \textbf{\exam{}}}
    \vspace{0.7cm}

    \begin{flushright}
        \textbf{Time Limit:} 50 minutes
    \end{flushright}
    \vspace{-0.2cm}
    
    \fbox{\parbox{\textwidth}{\centering 
        \begin{enumerate}[leftmargin = 0.5cm, label = $\bullet$, rightmargin = 0.2cm]
    
            \item \textbf{Every non-trivial calculation and claim in the exam must be properly justified.}

            \item The use of calculators is not necessary and therefore not allowed.

            \item Digital devices such as mobile phones, tablets, and laptops, as well as internet access, are forbidden.
            
        \end{enumerate}
    }}
\end{center}
%--- Instructions ---%

%--- Questions ---%

\begin{questions}

% Question
\question[3]
    Consider the demand and supply functions
\[
p_D(q)=21-(q+1)^2,
\qquad
p_S(q)=e^{\,q+\ln(12)-2}.
\]

\begin{enumerate}[leftmargin=0.5cm,label=(\alph*),rightmargin=0.2cm]
  \item Prove that at $q^*=2$ there is a supply--demand equilibrium. What is the equilibrium price $p^*$?

  \item From the point of view of a monopoly, the revenue of selling $q$ items (each at $p_D(q)$ euros) is $R(q)=p_D(q)\,q$. What value of $q$ maximizes the revenue?

  \item What is the inverse of the supply function? That is, find the function $q_S(p)$ such that $(q_S\circ p_S)(q)=q$. \\ 
  (Recall that $q$ must be positive as it is a quantity)
\end{enumerate}
% Solution
\begin{solution}[0cm]
\begin{enumerate}[leftmargin=0.5cm,label=(\alph*),rightmargin=0.2cm]
  \item Equilibrium solves $p_D(q)=p_S(q)$. At $q=2$,
  \[
    p_D(2)=21-(2+1)^2=21-9=12,\qquad
    p_S(2)=e^{\,2+\ln(12)-2}=e^{\ln(12)}=12.
  \]
  Hence $q^*=2$ and $p^*=12$.
  \item $R(q)=q\,p_D(q)=q\bigl(21-(q+1)^2\bigr)=20q-2q^2-q^3$. Then
  \[
    R'(q)=20-4q-3q^2=-(3q+10)(q-2).
  \]
  Critical points: $q=-\tfrac{10}{3}$ (discard) and $q=2$. Since $R''(q)=-4-6q<0$ for $q>0$, $q=2$ maximizes revenue.
  \item From $p=e^{\,q+\ln(12)-2}=12\,e^{\,q-2}$,
  \[
    \ln(p)=\ln(12)+q-2 \;\Longrightarrow\; q_S(p)=2+\ln\!\Big(\frac{p}{12}\Big).
  \]
  Since $q>0$, the relevant domain for the inverse is $p\ge 12e^{-2}$.
\end{enumerate}
\end{solution}


% Question
\question[2]
    You have a credit card debt that increases $1\%$ every month. How long will it take for the debt to double?    
%% Solution
\begin{solution}[0cm]
    Solve $(1.01)^n=2$ for months $n$:
    \[
        n=\frac{\ln(2)}{\ln(1.01)}\approx \frac{0.6931}{0.009950}\approx 69.7\ \text{months} \approx 5.8\ \text{years}.
    \]
\end{solution}

% Question
\question[2]
    Consider the functions $f(x) = \frac{\sqrt{x+2} - 2}{\sqrt{x - 2}}$ and $g(x) = \frac{1}{(x - 2)^2} + x^{-1} + x$, as well as the piecewise defined function
    \begin{align*}
        h(x) =
        \begin{cases}
            f(x), & x > 2 \\
            0, & x = 2 \\
            g(x), & x < 2
        \end{cases}
    \end{align*}

    \begin{enumerate}[leftmargin = 0.5cm, label = (\alph*), rightmargin = 0.2cm]
        
        \item What are the domains of $f(x)$ and $g(x)$?

        \item Is $h(x)$ continuous from the left at $x = 2$? Is it continuous from the right?

        \item Find $\lim_{x\rightarrow-\infty} h(x)$.
        
    \end{enumerate}
    
%% Solution
\begin{solution}[0cm]
    \begin{enumerate}[leftmargin = 0.5cm, label = (\alph*), rightmargin = 0.2cm]
        \item $f$: need $x-2>0$ and $x+2\ge0\Rightarrow x>2$.  
              $g$: exclude $x=2$ and $x=0\Rightarrow(-\infty,0)\cup(0,2)\cup(2,\infty)$.
        \item Left of $2$, $h=g$ and $g(x)\to+\infty$ as $x\to2^-$, so not left-continuous.  
              Right of $2$, rationalize
              \[
              f(x)=\frac{\sqrt{x+2}-2}{\sqrt{x-2}}
              =\frac{x-2}{(\sqrt{x+2}+2)\sqrt{x-2}}
              =\frac{\sqrt{x-2}}{\sqrt{x+2}+2}\to 0,
              \]
              which equals $h(2)=0$; thus right-continuous.
        \item For $x\to-\infty$, $h(x)=g(x)$ and
              \[
              \frac{1}{(x-2)^2}\to0,\quad x^{-1}\to0,\quad x\to-\infty \;\Rightarrow\; \lim_{x\to-\infty}h(x)=-\infty.
              \]
    \end{enumerate}
\end{solution}

% Question
\question[3]
Let $x(t)$ be the cumulative number of units a company has produced by month $t$, and let $C(x)$ be the total cost (in thousands of euros) of producing $x$ units.
Define $c(t) = C(x(t))$.
\begin{enumerate}[leftmargin = 0.5cm, label = (\alph*), rightmargin = 0.2cm]
    \item Provide an interpretation of $c(t)$.
    \item Provide an approximation of $c(1.5)$, given that $C(100) = 5$, $x(1) = 100$, $C'(100) = 0.5$ and $x'(1) = 1$.
    \item Let $P(t) = r(t) - c(t)$ be the profit of the company at time $t$, where $r(t)$ represents the revenue at $t$. Assume that $P(t)$ reaches a maximum at $t^*$. 
    \begin{enumerate}[leftmargin = 0.5cm, label = (\roman*), rightmargin = 0.2cm]
        
        \item What relation should $r'(t^*)$ and $c'(t^*)$ satisfy?

        \item What is the sign of $P'(t)$ for values of $t$ near $t^*$?  
        
    \end{enumerate}
\end{enumerate}

% Solution
\begin{solution}[0cm]
\begin{enumerate}[leftmargin=0.5cm,label=(\alph*),rightmargin=0.2cm]
  \item $c(t)$ is the total cost (in thousands of euros) as a function of time.
  \item By the chain rule $c'(t)=C'(x(t))\,x'(t)$, so $c'(1)=C'(100)\,x'(1)=0.5$. A linear approximation at $t=1$ gives
  \[
    c(1.5)\approx c(1)+c'(1)\,(1.5-1)=5+0.5\times0.5=5.25.
  \]
  \item Let $P(t)=r(t)-c(t)$.
  \begin{enumerate}[leftmargin=0.5cm,label=(\roman*),rightmargin=0.2cm]
    \item At a maximum $t^*$, \(P'(t^*)=0\). Since \(P'(t)=r'(t)-c'(t)\), the condition is
    \[
      r'(t^*)=c'(t^*).
    \]
    This is the reason behind the economic principle
    \begin{center}
    ``profit maximization is achieved when marginal cost equals marginal revenue''.    
    \end{center}
    
    \item Near a local maximum, \(P'(t)\) changes sign from non-negative to non-positive:
    \[
      P'(t)\geq 0 \text{ for } t<t^* \text{ close to } t^*,\qquad
      P'(t)\leq 0 \text{ for } t>t^* \text{ close to } t^* .
    \]
  \end{enumerate}
\end{enumerate}
\end{solution}

\end{questions}	

\newpage
\begin{center}
You may find useful the following derivatives of basic functions and rules of derivatives
\end{center}

\begin{center}
    \textbf{Differentiation rules}
\end{center}
        
        \begin{enumerate}[leftmargin = 0.5cm, label = , rightmargin = 0.2cm]

        \item \textbf{Linearity:} $h(x) = a f(x) + b g(x) \Longrightarrow h'(x) = a f'(x) + b g'(x)$, for $a, b \in\R$

        \item \textbf{Product rule:} $h(x) = f(x)g(x) \Longrightarrow h'(x) = f'(x)g(x) + f(x)g'(x)$

        \item \textbf{Quotient rule:} $h(x) = f(x)/g(x) \text{ and } g(x) \neq 0 \Longrightarrow h'(x) = \frac{f'(x)g(x) - f(x)g'(x)}{g(x)^2}$

        \item \textbf{Chain rule:} $h(x) = g(f(x)) \Longrightarrow h'(x) = g'(f(x))f'(x)$

        \item \textbf{Inverse rule:} $h(x) = f^{-1}(y) \Longrightarrow h'(y) = 1/f'(f^{-1}(y))$
        
        \end{enumerate}

\begin{center}
    \textbf{Derivative of usual functions}
\end{center}

        \begin{enumerate}[leftmargin = 0.5cm, label = , rightmargin = 0.2cm]

        \item \textbf{Constant:} $f(x) = a,\; a \in\R \Longrightarrow f'(x) = 0$
        
        \item \textbf{Power:} $f(x) = x^a \Longrightarrow f'(x) = a x^{a-1}$
        

        \item \textbf{Exponential:} $f(x) = e^x \Longrightarrow f'(x) = e^x$

        \item \textbf{Logarithm:} $f(x) = \ln(x) \Longrightarrow f'(x) = 1/x$

        \item \textbf{Trigonometric:} $f(x) = \sin(x),\; g(x) = \cos(x) \Longrightarrow f'(x) = \cos(x),\; g'(x) = -\sin(x)$
        
        \end{enumerate}

\end{document}
